%%% SETUP %%%
\documentclass[11pt]{article}

\usepackage[margin=0.9in]{geometry}
\usepackage{graphicx}
\usepackage{float}
\usepackage{amsmath,amssymb}
\usepackage{booktabs}
\setlength{\parskip}{2pt}
\setlength{\intextsep}{6pt plus 1pt minus 1pt}
\setlength{\textfloatsep}{6pt plus 1pt minus 1pt}
\setlength{\abovecaptionskip}{2pt}
\setlength{\belowcaptionskip}{2pt}

\title{\textbf{Predicting Soil Organic Carbon from AlphaEarth Foundation
Embeddings: An Exploratory and Inferential Analysis}}
\author{Luke Robinson \\ STAT 5000}
\date{April 24, 2026}

\begin{document}
\maketitle
\hrule
\vspace{0.5em}

\begin{abstract}
\noindent We assess whether the 64-dimensional geospatial embeddings produced
by Google DeepMind's AlphaEarth Foundations model carry usable signal for
predicting Soil Organic Carbon (SOC). Using $n = 3000$ random points in the
contiguous United States, paired with OpenLandMap SOC and MODIS land-cover
labels, we apply three techniques from STAT~5000 ---
\textbf{exploratory data analysis}, \textbf{Welch's two-sample
hypothesis test}, and \textbf{nonparametric bootstrap confidence
intervals} --- and find that several embedding dimensions are strongly
correlated with $\log(1+\mathrm{SOC})$ ($|r|$ up to $0.61$, BCa 95\%
CI $[-0.628, -0.581]$), and that mean log-SOC differs significantly
between Croplands and Grasslands. A supplementary spatial-block
cross-validated Random Forest reaches $R^2 = 0.75$ out-of-fold,
suggesting the embeddings encode subsurface-relevant signal beyond
what any single dimension captures.
\end{abstract}

%%% INTRO AND BACKGROUND %%%
\section{Introduction and Background}
\subsection{Research question}
Soil Organic Carbon (SOC) is one of the largest and most dynamic terrestrial
carbon pools and is central to both climate-mitigation accounting and
agricultural productivity \cite{lal2004soc}. Mapping SOC at scale is hard:
ground samples are expensive, sparse, and biased toward accessible land,
while traditional satellite-driven approaches require careful preprocessing of
multi-sensor imagery. Google DeepMind's AlphaEarth Foundations
\cite{brown2025alphaearth} sidesteps that preprocessing by compressing
Sentinel-1, Sentinel-2, Landsat, and climate inputs into 64 dense
unit-norm embedding bands at 10\,m resolution.

This project asks: \emph{do AlphaEarth embeddings carry usable signal for
SOC, even though they were not trained on any soil targets, and even though
SOC labels are sparse and indirect?}

\subsection{Background}
AlphaEarth is published as an annual mosaic in Google Earth Engine
\cite{gorelick2017gee} and was trained as a self-supervised representation of
the Earth's surface; the 64 dimensions have no direct physical interpretation.
For SOC labels we use OpenLandMap \cite{hengl2017openlandmap}, a global,
gridded SOC product at 250\,m derived by machine learning from a global
compilation of legacy soil profiles. Because both predictors and labels are
themselves model outputs, this study is best read as a test of whether the
\emph{AlphaEarth latent space encodes the same surface signals that drive an
established SOC model}, not as a ground-truth validation.

\subsection{Data}
We sampled $n_{0} = 4500$ random points uniformly within a CONUS bounding
box, dropped points where SOC or any AlphaEarth band was masked, and kept
the first $n = 3000$ valid records. For each point we extracted:
\begin{itemize}
\item AlphaEarth Foundations 2024 annual embedding, 64 bands $A_{00}..A_{63}$.
\item OpenLandMap SOC at 0\,cm depth (g/kg).
\item MODIS IGBP land cover \cite{friedl2010modis}, 2022.
\end{itemize}
All extraction was performed in Earth Engine with a 250\,m sampling scale to
match the SOC raster.

%%% METHODS AND RESULTS %%%
\section{Methods and Results}

\subsection{Method 1 --- Exploratory Data Analysis}
SOC is heavily right-skewed (skewness $= 2.39$, median $= 3$\,g/kg, max
$= 41$\,g/kg). The transformation $y = \log(1+\mathrm{SOC})$ is approximately
symmetric (Figure~\ref{fig:eda1}, left; mean $1.58$, sd $0.63$); subsequent
analyses use this log-transformed response.

\begin{figure}[H]
\centering
\begin{minipage}[b]{0.49\linewidth}
\centering
\includegraphics[width=\linewidth]{figures/eda-soc-hist-1.png}
\end{minipage}\hfill
\begin{minipage}[b]{0.49\linewidth}
\centering
\includegraphics[width=\linewidth]{figures/eda-boxplot-1.png}
\end{minipage}
\caption{Left: SOC distribution before and after the $\log(1+\cdot)$
transform. Right: log-SOC by IGBP land-cover class (top eight by $n$).
Forest classes have the highest median; the Grasslands ($n_1{=}1052$)
vs.\ Croplands ($n_2{=}543$) contrast motivates the hypothesis test below.}
\label{fig:eda1}
\end{figure}

We then computed the Pearson correlation of every AlphaEarth dimension
with log-SOC (Figure~\ref{fig:corr}); six dimensions exceed $|r| = 0.43$
and the strongest is $A_{40}$ with $r = -0.606$. This is a non-trivial
univariate signal given that the embeddings were never trained on any soil
target.

\begin{figure}[H]
\centering
\includegraphics[width=0.78\linewidth]{figures/eda-corr-bar-1.png}
\caption{Pearson correlation between each AlphaEarth dimension and
$\log(1+\mathrm{SOC})$. Six dimensions exceed $|r| = 0.43$; $A_{40}$ is
the strongest at $r = -0.61$.}
\label{fig:corr}
\end{figure}

\subsection{Method 2 --- Hypothesis Testing (Welch's Two-Sample $t$-test)}
We test whether the mean log-SOC differs between Grasslands ($n_1 = 1052$)
and Croplands ($n_2 = 543$), the two largest IGBP groups in our sample.
Formally,
\[
H_0: \mu_{\text{Grass}} = \mu_{\text{Crop}}, \qquad
H_1: \mu_{\text{Grass}} \neq \mu_{\text{Crop}}.
\]

\textbf{Assumption checks.} Per-group QQ plots of log-SOC are nearly
linear (notebook Section~3.2), and with $n_1 + n_2 = 1595$ the Central
Limit Theorem secures approximate normality of the sample means even
where per-observation deviation occurs. Levene's test does not reject
equal variances ($F_{1,1593} = 2.26$, $p = 0.13$); we still use Welch's
$t$ \cite{welch1947t} as a defensive choice since group sizes are unequal.

The Welch test statistic is
\[
t = \frac{\bar{y}_{\text{Grass}} - \bar{y}_{\text{Crop}}}
        {\sqrt{s_{\text{Grass}}^2/n_1 + s_{\text{Crop}}^2/n_2}}
  = -6.24, \quad \text{df}_{\text{Welch}} = 1326, \quad p = 6.0 \times 10^{-10}.
\]
The 95\% confidence interval for the mean difference
$\mu_{\text{Grass}} - \mu_{\text{Crop}}$ is $[-0.176, -0.092]$ on the log
scale, and Cohen's $d = -0.31$ (small-to-medium effect). \emph{We reject
$H_0$.} Croplands have higher mean log-SOC than Grasslands by roughly
0.13 log units --- equivalent to about a 14\% higher SOC on the natural
scale. This is consistent with the geography of the sample: U.S.~grasslands
are concentrated in the semi-arid West (low productivity, low residue
inputs), while croplands cluster in the Midwestern Mollisol belt with
historically high SOC stocks.

\subsection{Method 3 --- Bootstrap Confidence Intervals}
The Welch CI is exact only under normality. We complement it with
nonparametric bootstrap intervals \cite{efron1979bootstrap} that make no
distributional assumption. We use the \texttt{boot} package
\cite{canty2024boot} with $B = 10{,}000$ resamples and BCa
(bias-corrected and accelerated) intervals.

\textbf{(a) Bootstrap CI for the strongest correlation.} Resampling the
$(A_{40}, \log\text{-SOC})$ pairs and recomputing $r$ gives
$\hat{r} = -0.606$, BCa 95\% CI $[-0.628,\ -0.581]$ --- a narrow interval
that excludes zero by a wide margin (Figure~\ref{fig:boot}, left).

\textbf{(b) Bootstrap CI for the group mean difference.} A stratified
resample (preserving group sizes) for
$\bar{y}_{\text{Grass}} - \bar{y}_{\text{Crop}}$ gives a BCa 95\% CI of
$[-0.176, -0.091]$ (Figure~\ref{fig:boot}, right), agreeing with Welch's
CI $[-0.176, -0.092]$ to three decimal places --- a useful sanity check
that the parametric assumptions of Section~2.2 hold here.

\begin{figure}[H]
\centering
\begin{minipage}[b]{0.49\linewidth}
\centering
\includegraphics[width=\linewidth]{figures/boot-cor-fig-1.png}
\end{minipage}\hfill
\begin{minipage}[b]{0.49\linewidth}
\centering
\includegraphics[width=\linewidth]{figures/boot-md-fig-1.png}
\end{minipage}
\caption{Bootstrap distributions ($B=10{,}000$). Left: Pearson correlation
between $A_{40}$ and log-SOC. Right: stratified mean difference between
Grasslands and Croplands. Black: point estimate; red dashed: BCa 95\% CI.}
\label{fig:boot}
\end{figure}

\subsection{Beyond the three: PCA and Random Forest}
The above three methods come from the STAT~5000 syllabus. Two
supplementary analyses --- \emph{not counted toward the three} --- are
included because they directly answer the prediction question raised in
the project proposal.

\textbf{Principal Component Analysis.} A scree plot of a centered PCA on
the 64 embedding dimensions shows that 24 PCs are required to explain 90\%
of the variance and 32 PCs for 95\%; PC1 alone captures 20.8\%. The
embedding is therefore moderately --- but not extremely --- redundant.

\textbf{Spatial-block cross-validated Random Forest.} We fit a Random
Forest \cite{breiman2001randomforests,wright2017ranger} regressing log-SOC
on the 64 AlphaEarth dimensions, with 5-fold cross-validation in which
folds are assigned over $1^\circ \times 1^\circ$ lat/lon blocks rather
than over individual points (so neighbouring points cannot leak between
train and test). Mean out-of-fold $R^2 = 0.752$ with RMSE $= 0.32$ log-units.
A naive random-fold CV gives $R^2 = 0.758$, only marginally higher --- this
small gap is reassuring evidence that the embedding-to-SOC signal generalises
spatially rather than memorising location.

\begin{figure}[H]
\centering
\begin{minipage}[b]{0.48\linewidth}
\includegraphics[width=\linewidth]{figures/rf-pred-fig-1.png}
\end{minipage}\hfill
\begin{minipage}[b]{0.48\linewidth}
\includegraphics[width=\linewidth]{figures/rf-importance-1.png}
\end{minipage}
\caption{Left: out-of-fold predicted vs.~observed log-SOC, 5-fold
spatial-block CV. Right: top-15 AlphaEarth dimensions by Random Forest
impurity importance --- $A_{40}$, $A_{43}$, $A_{05}$, and $A_{33}$
(also the highest by univariate $|r|$) dominate.}
\label{fig:rf}
\end{figure}

\section{Conclusions}
We answered three questions and one bonus question:
\begin{enumerate}
\item Several AlphaEarth dimensions carry strong univariate signal for
SOC (Method 1; $|r|$ up to $0.61$).
\item Mean log-SOC differs significantly between Grasslands and Croplands
(Method 2; Welch's $t$, $p < 10^{-9}$).
\item The mean-difference and correlation effects are robust to
distributional assumptions (Method 3; BCa bootstrap intervals exclude zero
and agree closely with the parametric CI).
\item A Random Forest using all 64 dimensions predicts log-SOC with
$R^2 \approx 0.75$ on spatially-held-out blocks --- a useful baseline for
soil mapping that requires no traditional remote-sensing preprocessing.
\end{enumerate}

\subsection{Limitations}
Our SOC ``ground truth'' is itself an interpolated machine-learning product
\cite{hengl2017openlandmap}, so this work measures alignment between two
models rather than against field-measured SOC. The 250\,m SOC raster is
also coarser than AlphaEarth's 10\,m, and there is a temporal mismatch
(OpenLandMap is built on legacy profiles, AlphaEarth is annual 2024).
Spatial-block CV mitigates --- but does not eliminate --- spatial
autocorrelation effects.

\subsection{Future research}
Validating against field measurements (e.g., USDA SSURGO point profiles or
ISRIC's WoSIS database) would convert this from a model-vs-model study into
a true accuracy assessment. Extending to multi-depth SOC, pH, and texture,
and benchmarking AlphaEarth against a Sentinel-2-only baseline, would
isolate the marginal value of the foundation-model embedding.

%%% BIBLIOGRAPHY %%%
\bibliographystyle{plain}
\bibliography{references}

\end{document}
